\documentclass[11pt,a4paper]{article}

\newcommand{\ReportShortTitle}{Machine-checked cone algebra}
\newcommand{\ReportHeaderRight}{Scientific report}
\usepackage{reports/ipnpcns-report}

\hypersetup{
  pdftitle={A Machine-Checked Analysis of a Cone-Algebra Model for Neural Population Computation},
  pdfauthor={IPNPCNS formalization project},
  pdfsubject={Scientific report on Lean formalization and verification},
  pdfkeywords={Lean, formal verification, convex cones, Moreau decomposition, neural populations, NNLS}
}

\begin{document}

\hypersetup{pageanchor=false}
\begin{titlepage}
  \thispagestyle{empty}
  \begin{tikzpicture}[remember picture,overlay]
    \fill[ReportNavy] (current page.north west)
      rectangle ([yshift=-91mm]current page.north east);
    \fill[ReportCyan] ([yshift=-91mm]current page.north west)
      rectangle ([yshift=-95mm]current page.north east);
    \fill[ReportPale] ([yshift=27mm]current page.south west)
      rectangle (current page.south east);
  \end{tikzpicture}

  \vspace*{14mm}
  {\color{white}\sffamily\bfseries\fontsize{27}{32}\selectfont
    A Machine-Checked Analysis\par
    of a Cone-Algebra Model\par
    for Neural Population Computation\par}
  \vspace{7mm}
  {\color{white}\sffamily\Large
    Scientific report on the mathematics, formalization, and verification\par}

  \vfill

  \begin{tcolorbox}[
    colback=white,
    colframe=ReportCyan,
    boxrule=0.8pt,
    arc=1.5mm,
    left=4mm,right=4mm,top=3mm,bottom=3mm
  ]
  \begin{tabularx}{\linewidth}{@{}L{39mm}Y@{}}
    \textbf{Formal target} &
      \emph{Information Processing by Neuron Populations in the Central
      Nervous System: A Theory of the Mathematical Structure of Data and
      Operations}, arXiv:2309.02332v3 \\
    \textbf{Proof environment} &
      Lean 4.31.0 and mathlib revision
      \texttt{fabf563a7c95a166b8d7b6efca11c8b4dc9d911f} \\
    \textbf{Verified corpus} &
      31 Lean modules, 8,367 lines, and 439 theorem or lemma declarations \\
    \textbf{Source conformance} &
      Audited against the author-supplied LaTeX archive and an independent
      60-page source build \\
    \textbf{Date} & August 2026
  \end{tabularx}
  \end{tcolorbox}

  \vspace{12mm}
  {\large\color{ReportNavy}\textbf{IPNPCNS formalization project}}\par
  \vspace{2mm}
  {\color{ReportGray}
    A technical account for readers with mathematical maturity but no
    assumed knowledge of Lean.}
\end{titlepage}

\hypersetup{pageanchor=true}
\pagenumbering{roman}

\begin{abstract}
The model under study represents ensembles of neural population messages by
closed convex cones in a real Hilbert space. Population transformations are
then described using cone sum, intersection, reflection, metric projection,
and a residual operation called conic rejection. This report gives a
self-contained account of a Lean formalization of that mathematical program.

The central exact results are machine-checked for arbitrary closed convex
cones in arbitrary real Hilbert spaces. They include Moreau decomposition,
the identity that conic projection is two successive conic rejections, and
three rejection-only formulas for cone intersection. The development also
verifies the corresponding finite-subspace and Moore--Penrose formulas,
deterministic near-isometry and sparse Gram bounds, finite nonnegative
least-squares optimality and convergence, a finite collateral-occupancy
model, an exact orthonormal wavelet example, a compact causal Volterra
filter, a quantified membrane-update remainder, and conditional
population-level realization theorems.

The verification has a precise boundary. It proves mathematical consequences
of explicit assumptions; it does not prove that biological neural
populations satisfy those assumptions. Stochastic learning, generic
Laplace-domain filter realization, random embedding probabilities, support
selection by biological dynamics, and recurrent-network stability remain
outside the compiled theorem set. An equation-by-equation audit against the
LaTeX source found no material discrepancy in the claims made by the Lean
project. A full build contains no admitted proofs or custom axioms, and an
axiom audit across the principal endpoints reports only the standard
foundational principles used by Lean and mathlib.
\end{abstract}

\begingroup
\small
\setlength{\parskip}{0pt}
\tableofcontents
\endgroup
\clearpage

\pagenumbering{arabic}

\section{Question, answer, and scope}
\label{sec:introduction}

The mathematical proposal examined here begins with a change of object. A
single activity vector records one message from one neural population.
An ensemble of related messages, however, is represented by the closed convex
cone generated by those vectors. The invariant is therefore not a point but a
geometric set. A population may transform the coordinates of individual
messages while retaining the algebraic relations among the cones that those
messages generate.

This proposal creates three logically different questions:

\begin{enumerate}[label=\arabic*.]
  \item Are the definitions coherent, and do the claimed mathematical
    identities follow from the stated hypotheses?
  \item Do the idealized population and learning interfaces imply the claimed
    cone operations?
  \item Do real neural systems satisfy those interfaces with biologically
    meaningful error levels?
\end{enumerate}

Lean can settle the first question and conditional versions of the second.
The third is empirical. No proof assistant can infer anatomical connectivity,
time-scale separation, or in-vivo convergence from pure mathematics.

\begin{keyresult}[Overall conclusion]
The deterministic mathematical core is realistic to formalize and has been
formally verified at substantial breadth. The headline cone identities hold
without finite-dimensionality or finite-generation assumptions. The
population interpretation is verified only as an implication from explicit
model premises, not as a theorem about the central nervous system.
\end{keyresult}

The report is organized by mathematical dependence rather than by the order
in which files were written. Subspaces provide a linear benchmark. Closed
cones then supply the central geometry. Approximate invariance, temporal
signals, learning, occupancy, and population realization form separate
layers around that core. The final sections explain exactly what the machine
check certifies and where it stops.

\subsection{Four kinds of formal status}

Throughout the report, claims are classified as follows.

\begin{table}[htbp]
\centering
\caption{Interpretation of verification status.}
\label{tab:status}
\begin{tabularx}{\linewidth}{@{}L{30mm}Y@{}}
\toprule
\textbf{Status} & \textbf{Meaning} \\
\midrule
\statusverified &
A Lean theorem derives the stated conclusion from assumptions no stronger
than those shown. A formally stronger theorem may cover the paper statement
as a specialization. \\
\statusconditional &
The mathematical implication is proved, but a regularity, sampling,
training, or biological premise is made explicit rather than derived. \\
\statusrepresented &
An object or relation is encoded as an interface so that later consequences
can be checked, but the interface itself is not established by the
development. \\
\statusoutside &
The item is empirical, cited background, underspecified stochastic theory,
or a deliberately excluded mathematical extension. No verification claim is
made for it. \\
\bottomrule
\end{tabularx}
\end{table}

\section{From population messages to Hilbert-space geometry}
\label{sec:model}

\subsection{Messages and invariants}

Let \(H\) be a real Hilbert space. In the static finite-dimensional model,
\(H=\mathbb{R}^n\) and a message is a vector of population activities. In the
temporal model,
\[
  H=L^2([0,T];\mathbb{R}^n),
\]
so each neuronal coordinate is itself a time-varying square-integrable
signal. A finite family of messages \(v_1,\ldots,v_m\) generates the cone
\[
  C=\Conic\{v_1,\ldots,v_m\}
   =\left\{\sum_{i=1}^{m}\lambda_i v_i:\lambda_i\geq 0\right\}.
\]
For an arbitrary set \(S\subseteq H\), the source uses the more general
definition
\[
  \Conic(S)=
  \overline{\left\{\sum_{i=1}^{k}\lambda_i x_i:
  k\in\mathbb{N},\ x_i\in S,\ \lambda_i\geq0\right\}},
\]
where closure is in the norm topology. The closure is mathematically
important: pointwise images and Minkowski sums of closed cones need not be
closed in unrestricted infinite-dimensional settings.

The formal development therefore separates two levels.

\begin{itemize}
  \item Neural representations are usually finitely generated. Their spans
    remain finite-dimensional even if the ambient signal space is not.
  \item The main cone identities are stated for arbitrary closed convex cones
    in arbitrary real Hilbert spaces.
\end{itemize}

This separation both matches the source and strengthens the central results:
the proofs do not depend on a chosen finite frame.

\subsection{The operations}

For a closed cone \(C\), write \(P_Cx\) for the unique metric projection of
\(x\) onto \(C\). We use \(C^\circ\) for the source's non-positive polar:
\[
  C^\circ=\{z\in H:\langle z,y\rangle\leq 0
  \text{ for every }y\in C\}.
\]
The source prints this object as \(\neg C\). The present report uses
\(C^\circ\) to avoid confusing unary polarity with the binary rejection
operation.

\begin{table}[htbp]
\centering
\caption{Closed-cone operations used by the model. The report writes
\(A\crej B\) for the paper's binary conic-rejection notation.}
\label{tab:operations}
\begin{tabularx}{\linewidth}{@{}L{26mm}L{34mm}Y@{}}
\toprule
\textbf{Operation} & \textbf{Notation} & \textbf{Definition} \\
\midrule
Polar & \(C^\circ\) &
All vectors pairing non-positively with every point of \(C\). \\
Closed sum & \(A+B\) &
\(\Conic\{a+b:a\in A,\ b\in B\}\). \\
Reflection & \(-C\) &
\(\{-c:c\in C\}\). \\
Conic projection & \(A\cproj B\) &
\(\Conic\{P_Ba:a\in A\}\). \\
Conic rejection & \(A\crej B\) &
\(\Conic\{a-P_Ba:a\in A\}=A\cproj B^\circ\). \\
Intersection & \(A\cap B\) &
Ordinary set intersection, again a closed convex cone. \\
\bottomrule
\end{tabularx}
\end{table}

The operation \(A\cproj B\) is not merely the pointwise image
\(\{P_Ba:a\in A\}\). The metric projection onto a cone is generally nonlinear,
and its image of another cone need not be convex. Taking the closed conic
hull is part of the definition.

\subsection{Where verification sits}

\Cref{fig:boundary} summarizes the logical boundary. Arrows mean
\emph{is used as a premise for}, not \emph{is proved by}. The formal checker
reaches the top two mathematical layers. The empirical layer remains a
scientific validation problem.

\begin{figure}[htbp]
\centering
\begin{tikzpicture}[
  node distance=6mm,
  layer/.style={
    rounded corners=1.5mm,
    draw=ReportBlue,
    line width=0.7pt,
    minimum width=0.88\linewidth,
    text width=0.82\linewidth,
    align=left,
    inner sep=3.3mm
  },
  arr/.style={-{Latex[length=2.4mm]},line width=0.8pt,draw=ReportGray}
]
  \node[layer,fill=ReportPale] (proof) {
    \textbf{\color{ReportNavy}Machine-checked mathematical consequences}\\
    Cone identities, projection bounds, NNLS optimality and convergence,
    finite occupancy expectations, and circuit rewrites.
  };
  \node[layer,fill=white,below=of proof] (interfaces) {
    \textbf{\color{ReportNavy}Explicit mathematical interfaces}\\
    Near-isometry coverage, Lipschitz drive, support-selected map,
    approximation error, independent sampling, and blocking control.
  };
  \node[layer,fill=ReportWarm,below=of interfaces,draw=ReportGold] (mechanism) {
    \textbf{\color{ReportInk}Mechanistic and anatomical interpretation}\\
    Synaptic signs, collateral structure, temporal filtering, learning,
    support selection, and network synchronization.
  };
  \node[layer,fill=white,below=of mechanism,draw=ReportRed] (evidence) {
    \textbf{\color{ReportRed}Empirical validity}\\
    Measurements establishing that actual neural populations satisfy the
    model with useful parameters.
  };
  \draw[arr] (evidence) -- (mechanism);
  \draw[arr] (mechanism) -- (interfaces);
  \draw[arr] (interfaces) -- (proof);
\end{tikzpicture}
\caption{Verification boundary. Lean checks consequences above the empirical
layer; it does not convert model interpretation into biological evidence.}
\label{fig:boundary}
\end{figure}

\section{What the machine check means}
\label{sec:lean}

\subsection{Proofs as checked objects}

Lean is an interactive theorem prover based on dependent type theory
\cite{demoura2015lean}. A proposition is represented by a type, and a proof is
a term inhabiting that type. Tactics help construct proof terms, but the small
kernel checks the final term independently of the tactic that produced it.
The mathematical library mathlib supplies definitions and previously checked
theorems for Hilbert spaces, convexity, measure theory, matrices, topology,
and finite combinatorics \cite{mathlib2020}.

For example, the formal headline theorem has the following readable shape:

\begin{lstlisting}[style=leanstyle,language={},caption={The public Lean endpoint for double rejection.}]
theorem doubleRejection (A B : ClosedCone H) :
  coneProjection A B =
    coneRejection A (coneRejection A B)
\end{lstlisting}

Here \(H\) is already declared to be a complete real inner-product space.
The colon separates the theorem name and parameters from the proposition
being proved. The declaration contains no algorithm for numerically computing
the projection. It is a proof of an extensional equality of closed cones.

\subsection{Classical and noncomputable mathematics}

Metric projection in an arbitrary Hilbert space is selected from an existence
theorem. The formal definition is therefore marked noncomputable. This is
standard classical mathematics: it permits reasoning about the unique
minimizer without claiming executable code for finding it. Uniqueness later
makes all observable results independent of the particular witness chosen.

The axiom audit reports only:

\begin{description}
  \item[Propositional extensionality.]
  Logically equivalent propositions may be identified.
  \item[Classical choice.]
  A witness may be selected from a proof that a suitable object exists.
  \item[Quotient soundness.]
  Equivalent representatives define the same quotient object.
\end{description}

These are standard foundations of Lean and mathlib. The development adds no
custom axiom and contains no admitted proof placeholder.

\subsection{What compilation does not mean}

A successful build establishes type correctness and proof checking for the
formal statements actually present. It does not establish that those
statements are faithful to the source merely because their names resemble
equation numbers. Faithfulness was addressed separately by reconstructing
the formulas and qualifying prose from the LaTeX source. It also does not
establish numerical stability of an implementation, because the central
metric projection is a mathematical object rather than extracted software.

\section{The linear benchmark: finite subspaces}
\label{sec:subspaces}

The paper introduces subspaces before cones. This layer is not dispensable
background: it fixes matrix orientation, projection signs, rejection
conventions, and comparison measures in a setting where familiar linear
algebra provides strong regression tests.

\subsection{Column spaces and the covariance rendering}

For a finite real matrix \(A\), the represented subspace is
\(\Col(A)\). The source derives
\[
  \Col(A)=\Col(AA^\mathsf{T}).
\]
The Lean proof is coordinate-free at its core: the range of a linear map
equals the range of its composition with its adjoint. The finite matrix
statement follows by translating matrices to Euclidean linear maps. It holds
for rectangular, zero, and rank-deficient matrices.

Subspace sum is represented by concatenating generators. Orthogonal
complement, projection, and rejection are represented by projector
composition. A two-dimensional exact counterexample confirms that sum does
not distribute over intersection.

\subsection{A Moore--Penrose inverse without a library black box}

The required version of mathlib did not provide a general bundled
Moore--Penrose matrix API. The formalization therefore constructs the
pseudoinverse mathematically.

For a finite-dimensional linear map \(A:E\to F\), restrict \(A\) to
\((\ker A)^\perp\). This restriction is a linear equivalence
\[
  A_{\mathrm{eff}}:(\ker A)^\perp \simeq \operatorname{range}(A).
\]
The inverse map is composed with orthogonal projection onto
\(\operatorname{range}(A)\), giving \(A^+:F\to E\). The development proves
all four Penrose equations:
\[
\begin{aligned}
  AA^+A&=A, &
  A^+AA^+&=A^+,\\
  (AA^+)^\ast&=AA^+, &
  (A^+A)^\ast&=A^+A,
\end{aligned}
\]
and proves uniqueness from those equations. Consequently, \(AA^+\) is the
orthogonal projector onto the column space and \(A^+A\) is the projector
onto the row space.

This construction closes a common informal gap. Formulas containing
\(A^+\) are not restricted to full-rank matrices unless that restriction is
explicitly stated.

\subsection{Subspace operations and comparison}

Let \(P,Q\) be closed subspaces. Intrinsic definitions use orthogonal
projectors rather than matrices:
\[
\begin{aligned}
  P\cproj Q
    &=\operatorname{range}(\pi_Q\circ\pi_P),\\
  P\crej Q
    &=P\cproj Q^\perp.
\end{aligned}
\]
The verified identities include
\[
  P\cproj Q=P\crej(P\crej Q)
\]
and
\[
  P\cap Q=(P+Q)\crej\bigl((Q\crej P)+(P\crej Q)\bigr).
\]
Matrix corollaries recover the source's pseudoinverse renderings.

For comparison, the development proves that the Frobenius pairing of
projector matrices equals a basis-independent overlap:
\[
  \langle \Pi_P,\Pi_Q\rangle_F
  =\tr(\Pi_P\Pi_Q)
  =\sum_{e\in\mathcal{B}_Q}\|\pi_Pe\|^2,
\]
for any orthonormal basis \(\mathcal{B}_Q\) of \(Q\). Normalizing this
quantity yields the directional and symmetric subspace measures in the
source.

\section{Closed cones, metric projection, and Moreau decomposition}
\label{sec:moreau}

\subsection{Metric projection}

Let \(C\) be a nonempty closed convex cone in a real Hilbert space \(H\).
The Hilbert projection theorem gives a unique point
\[
  P_Cx=\argmin_{y\in C}\|x-y\|.
\]
The formalization first selects a minimizer and proves membership and
minimality. It then proves the variational characterization:
\[
  y=P_Cx
  \quad\Longleftrightarrow\quad
  y\in C
  \text{ and }
  \langle x-y,w-y\rangle\leq0
  \quad(\forall w\in C).
\]
Because \(C\) is a cone, scaling test points turns this into the sharper
cone-polar characterization
\[
  y=P_Cx
  \quad\Longleftrightarrow\quad
  y\in C,\quad x-y\in C^\circ,\quad
  \langle y,x-y\rangle=0.
\]

\subsection{Moreau decomposition}

The central decomposition is Moreau's classical cone theorem
\cite{moreau1962}; the formal result below includes the uniqueness form used
throughout the later algebra.

\begin{theorem}[Moreau decomposition]
\label{thm:moreau}
For every closed convex cone \(C\subseteq H\) and every \(x\in H\),
\[
  x=P_Cx+P_{C^\circ}x,
  \qquad
  \langle P_Cx,P_{C^\circ}x\rangle=0.
\]
Conversely, every decomposition \(x=y+z\) with
\(y\in C\), \(z\in C^\circ\), and \(\langle y,z\rangle=0\) has
\(y=P_Cx\) and \(z=P_{C^\circ}x\).
\end{theorem}

The proof combines the projection characterization with bipolarity
\((C^\circ)^\circ=C\). A particularly useful corollary is
\[
  P_{C^\circ}x=x-P_Cx.
\]
This residual identity connects geometric projection to the model's
rejection operation.

\begin{figure}[htbp]
\centering
\begin{tikzpicture}[scale=1.05,>=Latex]
  \fill[ReportBlue!16] (0,0) -- (3.1,0) -- (0,2.35) -- cycle;
  \fill[ReportGray!18] (0,0) -- (-2.0,0) -- (0,-1.7) -- cycle;
  \draw[->,ReportGray] (-2.35,0) -- (3.45,0) node[right] {\(e_1\)};
  \draw[->,ReportGray] (0,-1.95) -- (0,2.65) node[above] {\(e_2\)};
  \draw[line width=1.0pt,ReportBlue] (0,0) -- (3.1,0);
  \draw[line width=1.0pt,ReportBlue] (0,0) -- (0,2.35);
  \draw[line width=1.0pt,ReportGray] (0,0) -- (-2.0,0);
  \draw[line width=1.0pt,ReportGray] (0,0) -- (0,-1.7);
  \node[ReportNavy] at (1.25,1.15) {\(C\)};
  \node[ReportGray] at (-0.85,-0.78) {\(C^\circ\)};
  \draw[->,line width=1.35pt,ReportNavy] (0,0) -- (2.2,-1.15)
    node[right] {\(x\)};
  \draw[->,line width=1.35pt,ReportBlue] (0,0) -- (2.2,0)
    node[above] {\(P_Cx\)};
  \draw[->,line width=1.35pt,ReportGold] (2.2,0) -- (2.2,-1.15)
    node[midway,right] {\(P_{C^\circ}x\)};
  \draw[ReportGray] (2.0,0) -- (2.0,-0.2) -- (2.2,-0.2);
\end{tikzpicture}
\caption{A two-dimensional Moreau decomposition. The cone component and
polar component are orthogonal and add to the original vector. The theorem
holds in arbitrary real Hilbert spaces.}
\label{fig:moreau}
\end{figure}

\subsection{Structural projection laws}

Two further source claims are verified globally:
\[
  P_C(\alpha x)=\alpha P_Cx
  \qquad(\alpha\geq0)
\]
and
\[
  P_Cx=P_C(P_{\Span(C)}x).
\]
The second identity says that components orthogonal to the real span of the
cone do not affect projection. For a finitely generated cone, the development
also proves that this span equals the span of its generators and is
finite-dimensional even when \(H\) is infinite-dimensional.

\section{The two headline cone identities}
\label{sec:headline}

\subsection{Projection by two rejections}

\begin{theorem}[Double rejection]
\label{thm:double}
For arbitrary closed convex cones \(A,B\) in a real Hilbert space,
\[
  A\cproj B=A\crej(A\crej B).
\]
\end{theorem}

The pointwise argument reveals why the identity is true. Fix \(x\in A\) and
write its Moreau decomposition relative to \(B\):
\[
  x=y+z,\qquad y=P_Bx,\qquad z=P_{B^\circ}x.
\]
Let \(R=A\crej B\). By definition, \(z\in R\), and
\(R\subseteq B^\circ\). Polarity reverses inclusion, so
\[
  B=(B^\circ)^\circ\subseteq R^\circ.
\]
Thus \(y\in R^\circ\), \(z\in R\), and \(y\perp z\). Moreau uniqueness,
now applied to the cone \(R^\circ\), gives
\[
  P_Bx=y=P_{R^\circ}x.
\]
This equality holds for every generator \(x\in A\). Taking the closed conic
hulls of the identical pointwise generator sets yields the theorem.

\begin{keyresult}[Why closure handling matters]
The proof does not claim that a pointwise projection image is already a
closed convex cone. It first proves equality of the pointwise generators and
then lifts that equality through the same closed-conic-hull constructor on
both sides.
\end{keyresult}

\subsection{Intersection by rejections and sums}

Define
\[
  R=(B\crej A)+(A\crej B),
  \qquad K=R^\circ.
\]
The source's three formulas are:
\[
\begin{aligned}
  A\cap B
    &=A\crej R,\\
    &=B\crej R,\\
    &=(A+B)\crej R.
\end{aligned}
\]

\begin{theorem}[Rejection-only intersection]
\label{thm:intersection}
All three formulas above hold for arbitrary closed convex cones in an
arbitrary real Hilbert space.
\end{theorem}

The longer proof has two geometric stages.

\begin{enumerate}
  \item Since every common point of \(A\) and \(B\) is polar to both directed
    rejection cones, \(A\cap B\subseteq K\). For \(q\in K\), four Moreau
    decompositions relative to \(A\), \(B\), and their directed residuals
    imply
    \[
      P_Aq=P_Bq=P_{A\cap B}q.
    \]
    The central estimate combines polar signs with a completed-square
    identity; nonpositivity and nonnegativity force both directed residuals
    to vanish.
  \item Project a point \(x\in A+B\) onto \(K\). The first stage decomposes
    that projection into a common component in \(A\cap B\) and a component
    polar to \(A+B\). Moreau uniqueness then shows that the projection onto
    \(K\) is exactly the common component. Hence
    \[
      \{P_Kx:x\in A+B\}=A\cap B.
    \]
    Restricting the source set to \(A\) or \(B\) gives the same image, and
    closed conic hulls yield the three cone equalities.
\end{enumerate}

The formal proof follows this architecture rather than replacing it with
finite-dimensional polyhedral computation. This is important: the theorem
retains the general quantification stated in the source.

\subsection{Algebraic consequences and circuits}

The exact cone layer also verifies:
\[
\begin{aligned}
  A\cproj B&\subseteq B,\\
  A\subseteq B&\Longleftrightarrow B^\circ\subseteq A^\circ,\\
  (A^\circ)^\circ&=A,\\
  A\cproj(A\cproj B)&=A\cproj B,\\
  A\cap B&=(A^\circ+B^\circ)^\circ.
\end{aligned}
\]
Reflection commutes with rejection:
\[
  -(Y\crej X)=(-Y)\crej(-X).
\]
These identities justify source-level circuits that implement projection
with two rejection primitives and intersection with three rejection
primitives, one standalone reflection, and closed cone sums.

\section{Comparison and approximate invariance}
\label{sec:approximation}

\subsection{Comparing two cones}

For a nonzero cone \(A\), let
\[
  U(A)=\{x\in A:\|x\|=1\}.
\]
The directed similarity from \(A\) to \(B\) is
\[
  s_\to(A,B)=\inf_{x\in U(A)}\|P_Bx\|,
\]
and the symmetric similarity is
\[
  s(A,B)=\min\{s_\to(A,B),s_\to(B,A)\}.
\]
The angle is \(\theta(A,B)=\arccos s(A,B)\). The formalization proves
\(0\leq s\leq1\), identity on nonzero cones, and the right-angle consequence
when one cone contains a direction polar to the other. The zero cone is
handled separately, as in the source.

A normalized finite frame gives a computable minimum \(\widehat{s}\). It is
always optimistic:
\[
  s(A,B)\leq\widehat{s}(A,B).
\]
If every unit direction in each cone lies within distance \(\gamma\) of a
normalized frame ray, nonexpansiveness of metric projection gives
\[
  0\leq\widehat{s}(A,B)-s(A,B)\leq\gamma.
\]
The coverage premise is named and visible. No sampled biological frame is
assumed to satisfy it automatically.

\subsection{Scaled near-isometry}

Let \(M:H\to H'\) be a bounded linear map. On a relevant set \(S\), a scaled
near-isometry has parameters \(\alpha>0\) and \(0\leq\rho<1\):
\[
  (1-\rho)\alpha\|u\|^2
  \leq\|Mu\|^2
  \leq(1+\rho)\alpha\|u\|^2,
  \qquad u\in S.
\]
When \(S\) is a real subspace, a balanced polarization argument proves
\[
  \left|
    \langle Mu,Mv\rangle-\alpha\langle u,v\rangle
  \right|
  \leq\rho\alpha\|u\|\,\|v\|.
\]
This preserves a positive or negative inner-product sign whenever its
unscaled magnitude exceeds \(\rho\|u\|\|v\|\). Boundary points with zero
margin are correctly excluded from a uniform sign guarantee.

\subsection{Projection error}

Let \(p=P_Bx\), and suppose \(q\in B\) has transformed-space optimality
\[
  \|M(x-q)\|\leq\|M(x-p)\|.
\]
If the near-isometry covers both residuals \(x-q\) and \(x-p\), then
\[
  \|x-q\|^2
  \leq\frac{1+\rho}{1-\rho}\|x-p\|^2.
\]
Metric-projection optimality gives the Pythagorean inequality
\[
  \|x-q\|^2\geq
  \|x-p\|^2+\|q-p\|^2.
\]
Combining them yields the verified error theorem
\[
  \|q-p\|
  \leq
  \sqrt{\frac{2\rho}{1-\rho}}\,\dist(x,B).
\]

The Lean endpoint takes feasibility, transformed optimality, and residual
coverage as explicit premises. It does not package the source's preceding
closed-image argument that constructs \(q\) as the preimage of
\(P_{MB}(Mx)\), nor the subsequent norm bound after reapplying \(M\).
The scalar calculation itself is fully checked.

\subsection{Sparse Gram bounds}

Let \(m_i\) be the columns of a transmission map, with common scale
\(\alpha>0\). Suppose
\[
  \left|\frac{\|m_i\|^2}{\alpha}-1\right|\leq\eta,
  \qquad
  \frac{|\langle m_i,m_j\rangle|}{\alpha}\leq\mu
  \quad(i\neq j).
\]
For coefficients supported on a finite set \(I\), define
\[
  M_Ix=\sum_{i\in I}x_im_i.
\]
The formal proof expands the squared norm through the active Gram matrix.
The diagonal error contributes at most
\(\eta\alpha\sum_i x_i^2\). For the off-diagonal part, the elementary
inequality
\[
  2|x_ix_j|\leq x_i^2+x_j^2
\]
counts exactly \(|I|-1\) partners for each active coordinate. Consequently,
\[
  \left|
    \|M_Ix\|^2-\alpha\|x\|_2^2
  \right|
  \leq
  \bigl(\eta+(|I|-1)\mu\bigr)\alpha\|x\|_2^2.
\]
For \(|I|\leq k\), this gives the source's restricted-isometry bound with
\(\rho_k=\eta+(k-1)\mu\). The proof is deterministic and uses no random
matrix theorem; its role here is the deterministic geometric part of the
sparse-recovery framework associated with restricted isometries
\cite{candes2005}.

\begin{boundarybox}[Boundary of the invariance layer]
The random Johnson--Lindenstrauss dimension estimate, exact equivariance of
the complete cone algebra, the branch-incidence specialization
\(\mu\leq r/d\), and the heterogeneous spatiotemporal path model are not
formal endpoints. The verified result begins with a bounded linear map and
explicit deterministic geometric premises.
\end{boundarybox}

\section{Signals, filters, and membrane dynamics}
\label{sec:signals}

\subsection{The recording-window signal space}

Lebesgue measure is restricted to the closed interval \([0,T]\). Scalar and
finite-vector signals are Bochner \(L^2\) equivalence classes over that
measure. A finite spatial matrix is lifted pointwise to a bounded linear map
on the signal space, with an operator-norm bound inherited from the spatial
map.

Finite nonnegative aggregation is verified in an arbitrary real normed
signal space. If all component outputs lie in a closed cone and all weights
are nonnegative, their finite weighted sum lies in the same cone. Temporal
filters are represented as bounded linear operators, and preservation of a
particular cone is a visible predicate rather than a consequence of
boundedness alone.

\subsection{A compact causal convolution model}

The general source discussion uses finite-window convolution with suitable
square-integrable kernels. Boundary conventions matter, so the development
constructs one nontrivial model explicitly: the causal prefix-integration
operator. The ambient Hilbert-space and compact-operator facts are standard
functional analysis \cite{conway1990}; the point at issue is the exact
finite-window realization.
\[
  (V_Tf)(t)=\int_0^t f(s)\,ds
  \quad\text{for almost every }t\in[0,T].
\]
For \(T\geq0\), the formalization proves
\[
  \|V_T\|\leq T.
\]
It represents the operator by an \(L^2\) kernel, approximates that kernel by
simple functions, converts each simple kernel into a finite-rank operator,
and obtains finite-rank approximation in operator norm. Compactness follows.

This model is identified with a causal unit-step impulse response under
causal truncation. Separate theorems state that it is not the periodic or
whole-line zero-extension convention. The generic Laplace transfer identity
remains an interface; the concrete construction does not pretend to prove a
general transform theorem.

\subsection{An exact finite wavelet dictionary}

The source's three wavelets are formalized in
\(L^2([0,T])\) with
\[
  T=0.20,\qquad
  h(t)=
  \begin{cases}
    \sin^2(\pi t/T),&0\leq t\leq T,\\
    0,&\text{otherwise},
  \end{cases}
  \qquad
  \gamma=\frac{4}{\sqrt{3T}}.
\]
The atoms are
\[
\begin{aligned}
  \phi_1(t)&=\gamma h(t)\cos(2\pi 30t),\\
  \phi_2(t)&=\gamma h(t)\sin(2\pi 30t),\\
  \phi_3(t)&=\gamma h(t)\cos(2\pi 60t).
\end{aligned}
\]
Exact Fourier integral calculations prove
\[
  \int_0^T\phi_i(t)\,dt=0,
  \qquad
  \langle\phi_i,\phi_j\rangle=\delta_{ij}.
\]
The unconstrained, nonnegative, and plasticity examples then reduce to exact
coordinate calculations. The verified residual and neuronal-output signs
match the source. Literal firing-rate baselines are kept separate from the
baseline-subtracted \(L^2\) signals used in projection.

\subsection{A quantified membrane update}

Let \(d(s)\) be the net excitation-minus-inhibition drive over a local step,
\(\tau>0\), and \(\Delta t\geq0\). The exact first-order low-pass response is
\[
  z(\Delta t)=e^{-\Delta t/\tau}z_0+
  \int_0^{\Delta t}
    \frac{1}{\tau}e^{-(\Delta t-s)/\tau}d(s)\,ds.
\]
With
\[
  \lambda=1-e^{-\Delta t/\tau},
\]
the integral of the kernel is exactly \(\lambda\). Adding and subtracting the
frozen drive \(d(0)\) yields
\[
  z(\Delta t)
  =(1-\lambda)z_0+\lambda d(0)+r.
\]
If \(d\) is \(L\)-Lipschitz on the step, the formalization constructs the
remainder and proves
\[
  \|r\|\leq\frac{L}{\tau}|\Delta t|^2.
\]
Thus the source's \(O(\Delta t^2)\) is not an opaque annotation. It has an
explicit witness and bound. The premise is deliberately not inferred for an
arbitrary \(L^2\) signal, because an \(L^2\) equivalence class need not have
pointwise Lipschitz regularity.

\section{Nonnegative least squares and learning}
\label{sec:nnls}

The optimization layer uses the classical nonnegative least-squares problem
\cite{lawson1995}, but verifies the particular fixed-point, optimality, and
convergence claims needed by the model instead of importing an algorithm as
a black box.

\subsection{Finite objective, gradient, and projection}

Let \(X\in\mathbb{R}^{n\times N}\), \(y\in\mathbb{R}^{N}\), and
\(w\in\mathbb{R}^n\). The finite batch objective is
\[
  F(w)=\frac12\|X^\mathsf{T}w-y\|_2^2,
  \qquad w\geq0.
\]
Its gradient is
\[
  g(w)=X(X^\mathsf{T}w-y).
\]
The formalization proves the exact quadratic expansion
\[
  F(w+h)=F(w)+\langle g(w),h\rangle+
  \frac12\|X^\mathsf{T}h\|_2^2,
\]
and proves that \(g(w)\) is the directional derivative in every coefficient
direction.

Projection onto the nonnegative orthant is componentwise positive part.
For \(\varepsilon>0\), define
\[
  T_\varepsilon(w)=
  P_{\mathbb{R}_+^n}(w-\varepsilon g(w)).
\]
The formal theorem states
\[
  T_\varepsilon(w)=w
  \quad\Longleftrightarrow\quad
  w\geq0,\quad g(w)\geq0,\quad
  w_i g_i(w)=0\ \text{for every }i.
\]
The conditions on the right are the complementary KKT conditions. Using the
quadratic expansion, the development proves both directions of
\[
  \text{KKT point}\quad\Longleftrightarrow\quad
  \text{global NNLS minimizer}.
\]

\subsection{Deterministic convergence}

Orthant projection is nonexpansive. Under lower and upper spectral bounds
\[
  \mu\|h\|_2^2
  \leq\|X^\mathsf{T}h\|_2^2,
  \qquad
  \|Xz\|_2^2\leq L\|z\|_2^2,
\]
with \(0<\mu\leq L\) and \(0<\varepsilon<2/L\), the unprojected error map has
contraction factor
\[
  q=\sqrt{1-\varepsilon(2-\varepsilon L)\mu}<1.
\]
The projected iteration inherits this contraction and converges
geometrically to the unique coefficient minimizer.

The formal result is stronger in the rank-deficient case than a simple
contraction proof. With only an upper spectral bound, a nonempty solution
set, \(L>0\), and \(0<\varepsilon<2/L\), it proves:

\begin{itemize}
  \item Fejer descent toward every minimizer;
  \item summability and vanishing of prediction error;
  \item convergence in finite coefficient space to some NNLS minimizer;
  \item asymptotic regularity of successive iterates; and
  \item equality of the predictions produced by all minimizers.
\end{itemize}

Coefficient uniqueness is recovered only when \(w\mapsto X^\mathsf{T}w\)
is injective. Kernel directions are not incorrectly collapsed.

\subsection{The stochastic boundary}

One application of the clipped online rule is exactly a projected
stochastic-gradient step when the output is the instantaneous residual.
Convergence of a stochastic sequence is a different theorem. It requires a
specified probability space or filtration, unbiased sampling, moment
bounds, and step-size conditions. The source states these as standard
qualifications; the Lean project does not invent the missing stochastic
model. General projected stochastic-gradient results illustrate the kind of
additional hypotheses that such a theorem needs \cite{geiersbach2019}.

Similarly, a constant small learning rate generally tracks a solution with
steady-state misadjustment rather than converging exactly. When the membrane
state has \(\lambda<1\), the transient product is an exact gradient only for
a correspondingly matched filtered objective. These limits are part of the
formal interpretation, not footnotes to be discarded.

\section{Finite collateral occupancy}
\label{sec:occupancy}

The occupancy calculation is formalized as a finite probability experiment,
implemented by exact uniform averages over finite sets. One active axon
chooses a \(p\)-element subset of \(n\) targets. A sample of \(m\) axons is a
function assigning one such subset to each axon; the Cartesian product makes
the choices independent.

For a fixed target, the miss probability for one axon is \(1-p/n\), hence
the miss probability for all \(m\) axons is \((1-p/n)^m\). Summing target
indicators gives:

\begin{theorem}[Exact occupied-target expectation]
For \(0<n\) and \(p\leq n\), the expected number of occupied targets is
\[
  q=n\left(1-\left(1-\frac{p}{n}\right)^m\right).
\]
\end{theorem}

Zero targets, zero axons, zero collaterals, and invalid \(p>n\) cases have
separate theorems, so division and empty sample spaces are not hidden inside
the main statement.

The source replaces the finite miss probability by an exponential. Define
\[
  q_{\exp}=n\left(1-e^{-mp/n}\right).
\]
The formalization proves the signed and absolute bounds
\[
  q_{\exp}\leq q,
  \qquad
  |q-q_{\exp}|\leq\frac{mp^2}{n}.
\]
This replaces an informal approximation sign by a uniform finite inequality.

For two messages sampled independently from the same product experiment,
targetwise indicator independence yields the exact overlap formula
\[
  \mathbb{E}|O_1\cap O_2|=\frac{q^2}{n}.
\]
The theorem is exact for that product law. Its biological use still depends
on the modeling approximation that distinct messages have the required
independence and uniform marginals.

\section{Population-level realization}
\label{sec:population}

\subsection{Active faces}

For a finite generator family \(g_i\), let \(C_I\) be the closed cone
generated by indices \(I\), and let \(J\subseteq I\) be a candidate active
set. Orthogonal projection onto the span of \(\{g_j:j\in J\}\) is a candidate
piece of the nonlinear cone projection.

The formal certificate requires:

\begin{enumerate}
  \item the span projection has a nonnegative expansion in the selected
    generators \(g_j\);
  \item the residual pairs non-positively with every ambient generator
    \(g_i\), \(i\in I\).
\end{enumerate}

These are finite feasibility and polar inequalities. They show that the
candidate point lies in \(C_I\), the residual lies in \(C_I^\circ\), and the
two are orthogonal. Moreau uniqueness then proves
\[
  P_{C_I}x=P_{\Span\{g_j:j\in J\}}x.
\]
This is a constructive sufficient form of the source's local-face equation.
Non-strict inequalities deliberately admit boundary overlap. A separate
theorem proves that any two valid certificates at the same point return the
same cone projection.

The development does not prove that these regions are exactly the intrinsic
relative interiors of all polyhedral faces, nor that a biological threshold
mechanism selects them. Those are stronger geometric and empirical claims.

\subsection{Approximate population maps}

A population region records:

\begin{itemize}
  \item a source cone and target cone;
  \item a modeled message family;
  \item one active region;
  \item the corresponding face-span projection; and
  \item a support-selected population map \(T\).
\end{itemize}

The source's learned approximation premise is exposed as
\[
  \|T(u)-P_{\Span(F)}u\|\leq\epsilon\|u\|
\]
on the relevant intersection of source cone, message family, and active
region. Since active-face certification gives
\(P_{\Span(F)}u=P_Qu\), Lean proves
\[
  \|T(u)-P_Qu\|\leq\epsilon\|u\|
\]
and the identical residual error
\[
  \|(u-T(u))-(u-P_Qu)\|\leq\epsilon\|u\|.
\]
In the exact zero-error case, if the complete source cone is covered, the
closed conic hulls of outputs and residuals are exactly
\[
  \Conic(T(P))=P\cproj Q,
  \qquad
  \Conic\{u-T(u):u\in P\}=P\crej Q.
\]

\begin{boundarybox}[What equation (204) means formally]
The learned approximation bound is a premise. Lean proves its consequences
and prevents it from being hidden inside a definition asserting that a
neuron implements projection. It does not derive the numerical value of
\(\epsilon\) from anatomy, training data, or synaptic dynamics.
\end{boundarybox}

\subsection{Primitive circuits and conditionals}

The ideal output cone of one population is the reflected rejection
\[
  \mathcal{P}(P,Q)=-(P\crej Q).
\]
Using reflection equivariance, double rejection, and the intersection
identity, the development verifies the exact projection and intersection
circuits described in the source.

The conditional circuit has two explicit premises:

\begin{itemize}
  \item the control signal is inactive exactly when the condition cone is the
    zero cone;
  \item sufficiently strong synchronized inhibition passes the signal when
    control is inactive and blocks it otherwise.
\end{itemize}

Under those premises, the output is
\[
  \begin{cases}
    -(P\crej Q),&C=\{0\},\\
    \{0\},&C\neq\{0\}.
  \end{cases}
\]
No cone-emptiness detector, synchronization mechanism, or inhibitory
dynamics is smuggled into the proof.

\section{Verification evidence and traceability}
\label{sec:evidence}

\subsection{Authoritative-source audit}

The distributed paper was not the sole transcription source. The final
formal state was audited against the author-supplied archive
\texttt{arXiv-2309.02332v3.tar.gz}, whose SHA-256 digest is
\begin{center}
\small\ttfamily
a85f550647583374b73549e34cff596d0bd0c7991ea5586583f6cbdb068e1305.
\end{center}

The archive contains one top-level TeX source, an embedded bibliography, and
twelve PDF figures. Three independent \texttt{pdflatex} passes produced a
60-page document with no unresolved citation or cross-reference. A
layout-preserving text diff against the supplied arXiv PDF has one hunk:
the arXiv identifier and its induced title-page positioning. The article
body is identical.

The audit reconstructed formula bodies, labels, macro meanings, and nearby
qualifications through equation (210). It found no material mathematical
discrepancy in the claims made by the Lean project. One module-header section
reference was corrected, and the transmission-invariance limitations were
made more explicit in the project documentation.

\subsection{Build, gaps, and axioms}

The complete build checks 2,960 build jobs. The project contains 31 Lean
modules and 8,367 lines of Lean source. A direct scan finds no
\texttt{sorry}, custom \texttt{axiom}, or \texttt{unsafe} declaration in
project modules.

An explicit \texttt{\#print axioms} audit covers 38 public endpoints across
subspaces, cones, approximation, optimization, signals, wavelets,
population realization, circuits, and probability. Every endpoint depends
only on
\[
  \{\texttt{propext},\ \texttt{Classical.choice},\ \texttt{Quot.sound}\}.
\]
The job count and line count are reproducibility facts, not measures of
mathematical importance. The substantive evidence is the checked statement
of each endpoint and its source traceability.

\subsection{Coverage summary}

\begin{longtable}{@{}L{29mm}L{26mm}L{78mm}@{}}
\caption{Scientific coverage of the formal development.}
\label{tab:coverage}\\
\toprule
\textbf{Layer} & \textbf{Status} & \textbf{Verified content and boundary} \\
\midrule
\endfirsthead
\toprule
\textbf{Layer} & \textbf{Status} & \textbf{Verified content and boundary} \\
\midrule
\endhead
\bottomrule
\endfoot
Finite subspaces &
\statusverified &
Column-space identities, Moore--Penrose construction and uniqueness,
projection/rejection/intersection renderings, Frobenius and trace
comparisons, and non-distributivity. Gleason's theorem is cited background. \\

Closed cone core &
\statusverified &
Closed conic hull, non-positive polar, metric projection, Moreau
decomposition, closed sum, projection, rejection, reflection, exact cone
laws, double rejection, and intersection. \\

Cone comparison &
\statusverified &
Inclusion by zero rejection, exact worst-case similarity and angle, finite
frame proxy, and the \(\gamma\)-coverage error theorem. \\

Approximate invariance &
\statusconditional &
Near-isometry, polarization, angular margin, scalar projection error, and
sparse Gram/RIP bounds. Random embeddings, exact full equivariance, and the
anatomical branch specialization are outside the endpoints. \\

Temporal signals &
\statusconditional &
Recording-window \(L^2\), spatial lifting, finite aggregation, temporal
operator interfaces, compact causal unit-step Volterra filtering, exact
wavelets, and the Lipschitz membrane remainder. Generic Laplace and kernel
theory remain interfaces or outside scope. \\

NNLS and learning &
\statusverified &
Finite objective and gradient, projected fixed points, KKT equivalence,
full-rank contraction, rank-deficient convergence, prediction uniqueness,
and exact wavelet instances. General stochastic convergence is outside
scope. \\

Collateral occupancy &
\statusverified &
Exact finite expectation, explicit exponential error bound, and exact
two-message overlap under an independent product law. Biological
independence remains a model assumption. \\

Population realization &
\statusconditional &
Constructive sufficient active-face regions, consequences of the learned
approximation bound, exact hulls at zero error, primitive circuits, and the
abstract conditional. Training, support selection, gating dynamics, and
empirical error values are not derived. \\
\end{longtable}

\section{Scientific interpretation and limitations}
\label{sec:limits}

\subsection{What has been established}

The formalization establishes that the cone algebra is internally coherent
under a precise topological convention and that its two longest new
identities are valid in the generality stated by the source. It also
establishes that several surrounding quantitative claims can be made exact:
the occupancy experiment has a closed expectation formula, its exponential
replacement has an explicit finite error, the wavelet dictionary is exactly
orthonormal, the membrane remainder has a quadratic bound under Lipschitz
drive, and deterministic projected NNLS converges even when coefficient
minimizers are not unique.

The result is more than a transcription check. Important hidden assumptions
were either removed or exposed:

\begin{itemize}
  \item pseudoinverse formulas cover rank-deficient matrices;
  \item conic images use closed hulls explicitly;
  \item nonzero hypotheses appear in similarity measures;
  \item residual coverage appears in approximate-invariance theorems;
  \item temporal boundary conventions are distinguished;
  \item Lipschitz regularity is not inferred from \(L^2\);
  \item stochastic and deterministic convergence are separated; and
  \item population approximation and biological control remain premises.
\end{itemize}

\subsection{What has not been established}

The development does not show that a mammalian neural population represents
a cognitive concept by a convex cone. It does not validate the mechanistic
single-neuron model, estimate \(\rho\), \(\epsilon\), \(\gamma\), or
Lipschitz constants from data, or prove recurrent-network stability.
It does not formalize every equation in the source, every cited theorem, or
every illustrative population example.

The formal claim is narrower and sharper:

\begin{quote}
If messages, maps, filters, learning problems, and population regions satisfy
the formal interfaces stated here, then the verified cone-algebra,
approximation, optimization, and occupancy conclusions follow.
\end{quote}

This is the appropriate role of theorem proving in a mechanistic modeling
program. It isolates logical validity from empirical adequacy, allowing
future data or richer models to target named premises rather than an
informal chain of reasoning.

\subsection{Consequences for future work}

The most valuable extensions are not simply more lines of formal proof.
They are bridges that reduce scientifically important premises:

\begin{itemize}
  \item a concrete transform theorem connecting biological impulse responses
    to the Laplace-domain interface;
  \item a probability model for stochastic learning with explicit filtration
    and moment assumptions;
  \item an intrinsic polyhedral-face theory connecting certificates to
    relative-interior partitions and threshold support selection;
  \item certified data pipelines that turn measured connectivity and
    responses into bounds on \(\rho\), \(\mu\), \(\eta\), and \(\epsilon\);
    and
  \item network-level stability theorems for recurrent compositions of the
    verified population primitives.
\end{itemize}

Each extension would change the scientific boundary, not merely expand the
formal library.

\section{Conclusion}

The cone-algebra model has a substantial deterministic core that is both
mathematically meaningful and amenable to machine verification. The formal
development confirms Moreau-based projection and rejection, the double
rejection identity, the rejection-only intersection formulas, and a broad
supporting theory spanning subspaces, approximation, signals, optimization,
finite probability, and conditional population realization.

The strongest conclusion is therefore conditional but rigorous. Within the
explicit model class, neural population messages can be represented and
manipulated by an algebra of closed convex cones, and the central algebraic
claims are logically correct in arbitrary real Hilbert spaces. Lean checks
that conclusion without admitted proofs or project-specific axioms. Whether
the model describes biological computation with useful fidelity remains a
question for empirical science and for future formal bridges from data to
the named assumptions.

\clearpage
\appendix

\section{Source equations and principal Lean endpoints}
\label{app:traceability}

The table below is a compact index. Internal helper lemmas are omitted. Full
qualification and negative coverage are recorded in the source-conformance
audit.

\small
\begin{longtable}{@{}L{18mm}L{45mm}L{70mm}@{}}
\caption{Selected source-to-Lean traceability.}
\label{tab:endpoints}\\
\toprule
\textbf{Source} & \textbf{Mathematical content} & \textbf{Principal endpoint or artifact} \\
\midrule
\endfirsthead
\toprule
\textbf{Source} & \textbf{Mathematical content} & \textbf{Principal endpoint or artifact} \\
\midrule
\endhead
\bottomrule
\endfoot
(4)--(11) &
Column spaces and \(\Col(A)=\Col(AA^\mathsf{T})\) &
\nolinkurl{matrixColumnSpace_mul_transpose} \\

(12)--(23) &
Subspace complement, projection, rejection, and intersection &
\nolinkurl{double_subspaceRejection_eq_projection};
\nolinkurl{intersection_eq_sum_reject_mutualRejections} \\

(24), (27)--(29) &
Frobenius, trace, directional overlap, and cosine &
\nolinkurl{subspaceOverlap_eq_frobeniusInner};
\nolinkurl{symmetricSubspaceCosine_eq_frobenius} \\

(30) &
Occupancy expectation, exponential replacement, and independent overlap &
\nolinkurl{expectedOccupied_eq_exactFormula};
\nolinkurl{expectedOccupied_exponential_sq_error};
\nolinkurl{expectedOccupiedOverlap_eq_sq_div} \\

(31) &
Nonnegative finite signal aggregation &
\nolinkurl{aggregateSignal_mem_cone} \\

(33) &
Low-pass step and quadratic remainder &
\nolinkurl{exactMembraneUpdate_eq_updateWithRemainder} \\

(36), (39)--(41) &
Finite NNLS, gradient, fixed point, KKT, and global optimality &
\nolinkurl{nnlsLoss_hasDerivAt_line};
\nolinkurl{nnlsMinimizer_iff_complementaryKKT} \\

After (41) &
Full-rank and rank-deficient deterministic convergence &
\nolinkurl{projectedBatchIterate_tendsto_of_designSpectralBounds};
\nolinkurl{projectedBatchIterate_tendsto_some_nnlsMinimizer_of_upperBound_of_exists} \\

(62), (65)--(71) &
Closed conic hull, metric projection, polar, and Moreau decomposition &
\nolinkurl{conicHull}; \nolinkurl{metricProjection_eq_iff};
\nolinkurl{moreau_unique} \\

(72)--(80) &
Closed sum, conic projection, rejection, intersection, and reflection &
\nolinkurl{coneSum}; \nolinkurl{coneProjection}; \nolinkurl{coneRejection};
\nolinkurl{coneReflection} \\

(87)--(92) &
Exact and finite-frame cone comparison &
\nolinkurl{frameConeSimilarity_error92} \\

(93)--(97) &
Exact cone laws &
\nolinkurl{coneProjection_absorb};
\nolinkurl{inf_eq_polar_coneSum_polar} \\

(98)--(108) &
Projection by double rejection &
\nolinkurl{metricProjection_polar_rejection_eq};
\nolinkurl{doubleRejection} \\

(109)--(157) &
Intersection by rejections and sums &
\nolinkurl{intersectionByRejections} \\

(161)--(163) &
Scaled near-isometry, polarization, and sign margin &
\nolinkurl{scaledInnerDistortionOn_of_scaledNearIsometryOn_submodule};
\nolinkurl{inner_map_pos_of_margin};
\nolinkurl{inner_map_neg_of_margin} \\

(165)--(167) &
Projection-error calculation &
\nolinkurl{projection_error_of_scaledNearIsometryOn} \\

(170)--(171) &
Gain/coherence and sparse Gram bound &
\nolinkurl{sparseGramRIP171_of_card_le} \\

(176) &
Ideal population primitive &
\nolinkurl{populationPrimitive} \\

(195)--(196) &
Span reduction and nonnegative homogeneity &
\nolinkurl{metricProjection_starProjection_eq};
\nolinkurl{metricProjection_smul} \\

(199) &
Constructive active-face projection &
\nolinkurl{metricProjection_finiteGeneratedCone_eq_selectedFaceProjection} \\

(204), (206)--(208) &
Population approximation premise and verified consequences &
\nolinkurl{PopulationRegion.Approximation204};
\nolinkurl{PopulationRegion.projection_error206};
\nolinkurl{exact_population_hulls} \\

(209)--(210) &
Reflection/rejection and conditional circuit &
\nolinkurl{coneReflection_rejection};
\nolinkurl{controlledConditional210} \\
\end{longtable}
\normalsize

\section{Notation and assumption checklist}
\label{app:notation}

\begin{tabularx}{\linewidth}{@{}L{35mm}Y@{}}
\toprule
\textbf{Symbol or phrase} & \textbf{Meaning in this report} \\
\midrule
\(H\) &
An arbitrary real Hilbert space unless a finite Euclidean or recording-window
space is stated. \\
\(C^\circ\) &
The non-positive polar \(\{z:\langle z,c\rangle\leq0\ \forall c\in C\}\);
the source prints \(\neg C\). \\
\(\Conic(S)\) &
The norm-closed conic hull, not merely finite conic combinations. \\
\(P_Cx\) &
The unique metric projection onto a closed convex cone; a noncomputable
mathematical choice, not extracted numerical code. \\
\(A\cproj B\) &
Closed conic hull of pointwise projections of \(A\) onto \(B\). \\
\(A\crej B\) &
Closed conic hull of residuals \(a-P_Ba\); the paper's binary rejection. \\
Finite cone &
Finitely generated cone. The headline theorems need only closed convexity. \\
Near-isometry &
A deterministic two-sided squared-norm bound on every relevant intermediate
vector, not only on observed frame vectors. \\
Verified population map &
A theorem conditional on a support-selected map and approximation premise,
not an empirical claim about a recorded population. \\
\bottomrule
\end{tabularx}

\clearpage
\begin{thebibliography}{9}

\bibitem{nilsson2026}
M. N. P. Nilsson.
\newblock Information processing by neuron populations in the central nervous
system: A theory of the mathematical structure of data and operations.
\newblock arXiv:2309.02332v3, 2026.

\bibitem{moreau1962}
J.-J. Moreau.
\newblock Decomposition orthogonale d'un espace hilbertien selon deux cones
mutuellement polaires.
\newblock \emph{C. R. Acad. Sci. Paris}, 255:238--240, 1962.

\bibitem{lawson1995}
C. L. Lawson and R. J. Hanson.
\newblock \emph{Solving Least Squares Problems}.
\newblock SIAM, Philadelphia, 1995.

\bibitem{conway1990}
J. B. Conway.
\newblock \emph{A Course in Functional Analysis}.
\newblock Springer-Verlag, second edition, 1990.

\bibitem{candes2005}
E. J. Candes and T. Tao.
\newblock Decoding by linear programming.
\newblock \emph{IEEE Transactions on Information Theory},
51(12):4203--4215, 2005.

\bibitem{geiersbach2019}
C. Geiersbach and G. C. Pflug.
\newblock Projected stochastic gradients for convex constrained problems in
Hilbert spaces.
\newblock \emph{SIAM Journal on Optimization}, 29(3):2079--2099, 2019.

\bibitem{demoura2015lean}
L. de Moura, S. Kong, J. Avigad, F. van Doorn, and J. von Raumer.
\newblock The Lean theorem prover (system description).
\newblock In \emph{Automated Deduction -- CADE-25}, pages 378--388.
Springer, 2015.

\bibitem{mathlib2020}
The mathlib Community.
\newblock The Lean mathematical library.
\newblock In \emph{Proceedings of the 9th ACM SIGPLAN International
Conference on Certified Programs and Proofs}, pages 367--381, 2020.

\end{thebibliography}

\end{document}
